\chapter{Phimosis}

\section*{Definition}
Phimosis is the inability to retract the foreskin over the glans penis due to a narrowed preputial ring. It is physiologic in young boys (resolving in most by age 5--7) and pathological when it persists, is acquired, or causes symptoms.

\section*{What Happens in Your Body}
Physiologic phimosis results from natural adhesions between the inner foreskin and glans that dissolve over time as the penis grows and intermittent erections occur. Pathological phimosis involves cicatrix (scarring) of the preputial orifice, typically from recurrent balanitis, forceful retraction, or lichen sclerosus (balanitis xerotica obliterans), which produces a white, fibrotic ring.

\section*{Causes}
\begin{itemize}
  \item \textbf{Physiologic:} Normal developmental stage in uncircumcised children, resolving spontaneously in >90\% by puberty
  \item \textbf{Pathological:}
  \begin{itemize}
    \item Lichen sclerosus et atrophicus (balanitis xerotica obliterans) -- most common cause of pathological phimosis
    \item Recurrent balanitis or balanoposthitis
    \item Trauma from forced retraction, zipper injury, or catheterization
    \item Poor hygiene leading to smegma retention and chronic inflammation
    \item Diabetes (recurrent infections leading to scarring)
  \end{itemize}
\end{itemize}

\section*{When to See a Doctor}
See your doctor if:
\begin{itemize}
  \item Ballooning of the foreskin during urination (suggests significant obstruction)
  \item Urinary tract infections, recurrent balanitis, or difficulty with urination
  \item Painful erections or interference with sexual activity
  \item White, sclerotic, or scarred appearance of the preputial tip
\end{itemize}

\section*{Related Symptoms}
\begin{itemize}
  \item Balanitis, paraphimosis, painful urination, weak stream, foreskin ballooning, smegma retention, preputial pain, lichen sclerosus
\end{itemize}
\textbf{Important:} Forced retraction of the foreskin in a child should be avoided, as it can cause pain, fissuring, and pathological scarring.

\section*{Self-Care / Home Management}
\begin{itemize}
  \item For mild physiologic phimosis: gentle, gradual stretching over weeks to months using the thumbs to widen the preputial ring
  \item Apply emollient (petroleum jelly) to reduce friction during stretching
  \item Maintain good hygiene: clean visible areas only; do not attempt forced retraction
  \item Sitz baths to soothe inflammation
\end{itemize}

\section*{How Your Doctor Will Evaluate You}
\begin{itemize}
  \item History: age of onset, urinary symptoms, prior infections, forced retraction attempts
  \item Physical exam: inspect foreskin for scarring, balanitis, or lichen sclerosus (white plaque, loss of elasticity)
  \item Differentiate physiologic from pathological (scarring, pallor, inability regardless of relaxation)
  \item Assess for ballooning and post-void dribbling
\end{itemize}

\section*{Treatment Options}
\begin{itemize}
  \item First-line for pathological phimosis: topical corticosteroids (betamethasone 0.05\% cream or ointment) applied to the preputial ring for 4--8 weeks, with gentle stretching
  \item Success rate with topical steroids: 70--90\% for non-scarred phimosis
  \item Stretching exercises under steroid treatment
  \item Circumcision (surgical removal of the foreskin) for treatment failure or severe lichen sclerosus
  \item Preputioplasty (dorsal slit, Z-plasty) as a foreskin-preserving alternative
  \item Treat underlying balanitis or diabetes if present
\end{itemize}

\section*{References}
\begin{thebibliography}{99}
\bibitem{phimosis:phimosis1} Hayashi Y, Kojima Y, Mizuno K, et al. ``Phimosis: a review.'' Int J Urol. 2011;18(4):262-269.
\bibitem{phimosis:phimosis2} Srinath H. ``Phimosis in children.'' BMJ. 2020;370:m2775.
\bibitem{phimosis:phimosis3} McGregor TB, Pike JG, Leonard MP. ``Pathologic and physiologic phimosis: approach to the phimotic foreskin.'' Can Fam Physician. 2007;53(3):445-448.
\bibitem{phimosis:phimosis4} Esposito C, Centonze A, Alicchio F, et al. ``Topical steroid application versus circumcision in pediatric patients with phimosis.'' Pediatr Surg Int. 2008;24(2):187-191.
\bibitem{phimosis:phimosis5} Orsola A, Caffaratti J, Garat JM. ``Conservative treatment of phimosis in children using a topical steroid.'' Urology. 2000;56(2):307-310.
\bibitem{phimosis:phimosis6} Cuckow PM, Rix G, Mouriquand PDE. ``Preputial plasty: a good alternative to circumcision.'' J Pediatr Surg. 1994;29(4):561-563.
\end{thebibliography}
